% Conclusions are limited to the executed experiments (Section~\ref{sec:results}): the MLP
% study and the reduced-budget VGG-16-BN study.
We asked whether the learnable connection gates of an existing self-pruning MLP identify
prunable neurons. Every criterion used the same compact architectures, fine-tuning budget
and validation-based selection, over five seeds with 30-epoch sources. In the setting the
existing method ships with ($\lambda=0$) the answer is \emph{no}: its neuron score, the
mean incoming gate, gave lower accuracy than random neuron selection at every pruning ratio
and on every seed, by up to $4.4$\,pp after fine-tuning, and its scores were
anti-correlated with the measured damage of removing a neuron. The same result held on the
original 5-epoch checkpoints. This matches the analysis in Section~\ref{sec:analysis}:
behind batch normalization, the level of a neuron's incoming gates does not affect the
network function.

Three outcomes are constructive. First, reading the same gates on the outgoing side, which
BN does not cancel, recovers the accuracy of standard magnitude criteria. Second, any
positive gate penalty we tested removes the failure: gate-in then matches the other gate
scores and exceeds random selection (descriptively, three seeds). Even so, in the main study no gate score exceeded the conventional
criteria, and stronger sum penalties made the networks harder to prune for every
criterion. Third, the audit found and corrected a gate-folding error in the original
compact reconstruction. The results do \emph{not} show that gate-based criteria are better
than conventional ones, and they establish no ranking among the conventional criteria,
which were not tested against each other.

A reduced-budget VGG-16-BN study (20-epoch sources, three fine-tuning epochs, three seeds)
points in the same direction: gate-in fell below random at $40$ and $60\%$ channel removal
and lost to every other criterion on every seed at $60\%$. With three seeds, however, only
one comparison survives multiple-testing correction, so this is supporting rather than
independent evidence for convolutional networks. The main conclusions rest on one MLP
architecture on raw CIFAR-10 pixels, CPU-only measurements and three seeds for the
ablations. The full VGG-16-BN protocol (160-epoch sources, 20 fine-tuning epochs, five
seeds, gate-penalty ablation) on a GPU is the necessary next step before the findings can
be generalized.
